\def\arraystretch{1.6}
\begin{tabular}{@{}p{0.23\textwidth}p{0.33\textwidth}p{0.33\textwidth}@{}}
\toprule
\textbf{Methode}     & \textbf{Vorteile}            & \textbf{Nachteile}       \\
\midrule
Raumtemperatur                          & Die einfachste Methode. Am besten für Brot, das innerhalb eines Tages gegessen wird.
                                            Kruste bleibt typischerweise knusprig, wenn die Luftfeuchtigkeit nicht zu hoch ist.
                                          & Brot trocknet sehr schnell aus.\\

Raumtemperatur in luftdichtem Behälter    & Gut für bis zu eine Woche.
                                          & Brot muss getoastet werden, damit die Kruste wieder knusprig wird. Schimmelt schneller.\\

Kühlschrank                                    & Brot bleibt wochenlang gut. Kann ohne luftdichten Behälter etwas austrocknen.
                                          & Brot muss getoastet werden. Benötigt Kühlschrank und Energie.\\

Gefrierschrank                                   & Brot bleibt jahrelang gut.
                                          & Muss aufgetaut und dann getoastet werden. Benötigt Gefrierschrank und Energie.\\
\bottomrule
\end{tabular}
